\documentclass[11pt,a4paper]{article}
% Shared preamble for RuPhO English problem statements (match T1 2026 style).
% Use: \input{latex/rupho_en_preamble.tex} after \documentclass.
\usepackage[margin=1in]{geometry}
\usepackage{amsmath,amssymb}
\usepackage{graphicx}
\usepackage{float}
\usepackage{enumitem}
\usepackage{microtype}
\usepackage{caption}
\usepackage{tabularx}
\usepackage{hyperref}

\captionsetup{font=small,labelfont=bf,skip=6pt}

\setlist[itemize]{leftmargin=1.2cm}
\setlist[enumerate]{leftmargin=1.2cm}

% One outer box per subproblem: [id | statement | points] (no inner rules)
\newcommand{\problembox}[3]{%
  \par\addvspace{0.42em}%
  \noindent
  \setlength{\fboxsep}{7.5pt}%
  \setlength{\fboxrule}{0.95pt}%
  \fbox{%
    \begin{minipage}{\dimexpr\linewidth-2\fboxsep-2\fboxrule\relax}
    \begin{tabularx}{\linewidth}{@{}>{\raggedright\arraybackslash\bfseries}p{2.1em} @{\hspace{0.9em}} >{\raggedright\arraybackslash}X @{\hspace{0.9em}} >{\raggedleft\arraybackslash\bfseries}p{2.4em}@{}}
    #1 & #2 & #3
    \end{tabularx}
    \end{minipage}%
  }%
  \par\addvspace{0.32em}%
}


\title{\textbf{Flight to the Moon}}
\author{\normalsize Y26 (2026) --- English translation}
\date{}
\begin{document}
\maketitle

In this problem we will consider orbits in which you can fly to the Moon and return back. In the simplest problems, as a rule, a flight to the Moon is considered in an elliptical orbit with a major axis equal to $r_m$, where $r_m$ is the distance from the Earth to the Moon. Unfortunately, for a number of reasons, such an orbit is not the most optimal and safe. The main disadvantage of an elliptical orbit is the capture of the spacecraft in the gravitational field of the Moon. During this capture, the craft needs to enter an elliptical orbit around the Moon using maneuvering thrusters, and failure of the thrusters in such a situation could result in crashing onto the lunar surface. This safety problem is solved by a free-return trajectory. The moon acts as a gravitational slingshot and the trajectory appears to be a kind of figure eight. Such an orbit was used, for example, when sending the Soviet Luna-3 spacecraft to the Moon, and they also say that it was along it that the infamous Apollo 13 mission (which allegedly flew to the satellite of our beautiful planet) returned after an accident when approaching the Moon.

In the simplest problems, as a rule, a flight to the Moon is considered in an elliptical orbit with a major axis equal to $r_m$, where $r_m$ is the distance from the Earth to the Moon.

Unfortunately, for a number of reasons, such an orbit is not the most optimal and safe. The main disadvantage of an elliptical orbit is the capture of the spacecraft in the gravitational field of the Moon. During this capture, the craft needs to enter an elliptical orbit around the Moon using maneuvering thrusters, and failure of the thrusters in such a situation could result in crashing onto the lunar surface.

This safety problem is solved by a free-return trajectory. The moon acts as a gravitational slingshot and the trajectory appears to be a kind of figure eight. Such an orbit was used, for example, when sending the Soviet Luna-3 spacecraft to the Moon, and they also say that it was along it that the infamous Apollo 13 mission (which allegedly flew to the satellite of our beautiful planet) returned after an accident when approaching the Moon.

"Chronology" of the Apollo 13 mission.

\begin{figure}[H]
    \centering
    \includegraphics[width=0.72\textwidth]{T4_EN_assets/fig1.jpg}
    \caption{Figure 1.}
\end{figure}

Throughout the entire task, the following physical constants may be needed: gravitational constant $G = 6.67 \cdot 10^{-11}~\frac{м^3}{кг\cdot с^2}$ mass of the Earth $M_e = 6\cdot 10^{24}~кг$ mass of the Moon $M_m = \dfrac{M_e}{81}$ radius of the Earth $R_e = 6400~км$ radius of the Moon $R_m = 1700~км$ distance between the Earth and the Moon $r_m = 3.85 \cdot 10^8~м$ period of the Moon around the Earth $T_m = 27.4~суток$ orbital speed of the Moon $v_m = \dfrac{2\pi r_m}{T_m}$

Throughout the task, the following physical constants may be needed:

\section*{Part A. Takeoff from Earth (3.0 points)}

First, let's look at the process of taking off from Earth. A fairly well-known fact is that most of the rocket’s mass is made up of the first stage fuel, which is necessary to develop the first escape velocity and overcome the earth’s atmosphere. Accurate calculations of such takeoffs are extremely difficult. It is the complexity and the need to take into account an infinite number of factors that explains a considerable percentage of unsuccessful rocket launches. Nevertheless, some estimated values ​​can be obtained from fairly simple considerations. Let us consider the vertical takeoff of a rocket from the Earth along the $y$ axis (the $y$ axis is directed vertically upward, the initial coordinates and speed of the rocket are $0$). The initial mass of the rocket is $m_0$. The rocket is acted upon by gravity and the drag force $\vec{F} = -kv \vec{v}$. The rocket emits combustion products from the engine at a constant mass flow rate $\mu$ and speed $u$ relative to the rocket. Consider that $g$ varies little with height. Numerical values: $g = 9.8~\frac{м}{с^2}$ $\mu = 2700~\frac{кг}{с}$ $u = 4.5~\frac{км}{с}$ $k = 0.25~\frac{кг}{м}$ - value of the resistance force coefficient near the Earth’s surface $a = 6600 ~км$ $m_0 = 500 ~т$

First, let's look at the process of taking off from Earth. A fairly well-known fact is that most of the rocket’s mass is made up of the first stage fuel, which is necessary to develop the first escape velocity and overcome the earth’s atmosphere. Accurate calculations of such takeoffs are extremely difficult. It is the complexity and the need to take into account an infinite number of factors that explains a considerable percentage of unsuccessful rocket launches. Nevertheless, some estimated values ​​can be obtained from fairly simple considerations.

Let us consider the vertical takeoff of a rocket from the Earth along the $y$ axis (the $y$ axis is directed vertically upward, the initial coordinates and speed of the rocket are $0$). The initial mass of the rocket is $m_0$. The rocket is subject to gravity and drag force $\vec{F} = -kv \vec{v}$. The rocket emits combustion products from the engine at a constant mass flow rate $\mu$ and speed $u$ relative to the rocket. Consider that $g$ varies little with height.

Numerical values:

\problembox{A1}{Write the expressions for $\ddot y$ through $\dot{y}$, $k$, $m_0$, $\mu$, $u$, $g$ and $t$.}{0.50}

At the initial stage of takeoff, the rocket speed is low, so friction can be neglected. Let's find the limits within which this model works during vertical takeoff.

\problembox{A2}{Neglecting the drag force, find the speed of the rocket $v(m)$ and the coordinate $y(m)$ at the moment when its mass is $m$. Express your answer in terms of $m$, $g$, $u$, $\mu$ and $m_0$.}{1.00}

\problembox{A3}{Using the result of the previous paragraph and assuming that $k$ weakly depends on height, find numerically at what height $H_c$ the rocket will be at the moment when the drag force becomes equal to the force of gravity.}{0.80}

The height of the dense layers of the atmosphere is about $10$ km, so at high altitudes the drag force becomes negligible. Due to this fact, we will not take it into account in the next question.

\problembox{A4}{Considering a vertical take-off, numerically estimate the ratio $\beta_1$ of the initial to the final mass of the rocket, at which the rocket develops a speed $v_a$ (the speed of movement in a circular orbit of radius $a$).}{0.70}

In reality, the rocket is launched into orbit along a complex trajectory, so the resulting value does not exactly coincide with the actual value. In what follows, assume that the real value is $\beta_1 = 15$.

\section*{Part B. Free Return Trajectory (6.0 points)}

Let us move directly to the problem of considering the orbit of free return. First, let's find the area in which the gravitational attraction of the Moon affects more than the gravity of the Earth.

First, let's find the area in which the gravitational attraction of the Moon affects more than the gravity of the Earth.

\problembox{B1}{Consider a point on the Earth-Moon straight line between the Earth and the Moon, at which the gravitational forces of the Earth and the Moon are equal in magnitude. Find the distance $r_c$ from the Moon to this point. Express your answer in terms of $r_m$, $M_e$ and $M_m$.}{0.30}

\begin{figure}[H]
    \centering
    \includegraphics[width=0.72\textwidth]{T4_EN_assets/fig2.jpg}
    \caption{Figure 2.}
\end{figure}

In the future, IN ALL POINTS OF THE PROBLEM, we will call the sphere of radius $r_c$ around the Moon the zone of attraction of the Moon and assume that in it the change in the potential energy of the body is $\Delta U = \Delta U_m$ (i.e., the interaction with the Earth can be neglected), and in the rest of space $\Delta U = \Delta U_e$ (i.e., the interaction with the Moon can be neglected). Also, at all points, consider that the CO of the Moon is progressively moving. Let the ship initially be in a circular Earth orbit of radius $a$. At the right moment in time, it turns on the engines and accelerates to a speed of $v_0$ ($\vec v_0$ directed tangentially to the original circular orbit), such that it can reach the lunar gravitational zone. The minimum distance between the center of the Moon and the ship during movement is $b$. Right before entering the Moon's gravitational zone, the speed of the ship in the Earth's CO is $v_1$. Let us introduce the notation for the 2nd cosmic velocities for the corresponding orbits: $$v_\textbackslash\{\}text\{II\}^e = \textbackslash\{\}sqrt\{\textbackslash\{\}dfrac\{2GM_e\}\{a\}\}$$ $$v_\textbackslash\{\}text\{II\}^m = \textbackslash\{\}sqrt\{\textbackslash\{\}dfrac\{2GM_m\}\{b\}\}$$

Также во всех пунктах считайте, что СО Луны — поступательно движущаяся.

Пусть корабль изначально находится на круговой околоземной орбите радиуса $a$. В нужный момент времени он включает двигатели и разгоняется до скорости $v_0$ ($\vec v_0$ направлена по касательной к изначальной круговой орбите), такой, что может достигнуть зоны притяжения Луны. Минимальное расстояние между центром Луны и кораблем в процессе движения — $b$. Прямо перед входом в зону притяжения Луны скорость корабля в СО Земли равна $v_1$.

Введем обозначения для 2-ых коcmических скоростей для соответствующих орбит:

$$v_\textbackslash\{\}text\{II\}^e = \textbackslash\{\}sqrt\{\textbackslash\{\}dfrac\{2GM_e\}\{a\}\}$$

$$v_\textbackslash\{\}text\{II\}^m = \textbackslash\{\}sqrt\{\textbackslash\{\}dfrac\{2GM_m\}\{b\}\}$$

Пусть $\theta_1$ и $\theta_2$ углы показанные на рисунке. В общем виде орбита свободного возврата не обязана быть сиmmетричной, и траектория возврата может отличаться от траектории полета к Луне. Для простоты расчетов будем расcmатривать орбиту сиmmетричную относительно прямой Земля — Луна. Тогда угол $\theta_1 = \theta_2 = \theta$. Обозначим $\alpha$ угол между скоростью $v_1$ and direct Earth - Moon.

\begin{figure}[H]
    \centering
    \includegraphics[width=0.72\textwidth]{T4_EN_assets/fig3.jpg}
    \caption{Figure 3.}
\end{figure}

Further, to determine the required speed $v_1$, we will consider the orbit of the ship in the CO of the Moon in its zone of attraction. For convenience, we introduce the notations $s = \sin(\alpha + \theta)$, $s' = \sin \alpha$, $c = \cos \theta$.

\problembox{B2}{Considering the ship in the CO of the Moon, write down the expression for the specific angular momentum $L_0^m$ relative to the center of the Moon and the speed $v_2$ at the moment when the distance to the Moon is $r_c$. Express your answer in terms of $s$, $s'$, $c$, $v_1$, $v_m$ and $r_c$.}{0.50}

\problembox{B3}{By writing the ZSE in the Moon's CO for the position of the ship closest to the Moon and the position immediately after entering the attraction zone, obtain and solve the quadratic equation and obtain the exact expression for $v_1/v_m$. Express your answer in terms of $s$, $s'$, $c$, $b/r_c$ and $v_\text{II}^m/v_m$. Note. There is no need to select a root sign at this point. Do not forget that within the framework of the model used, the interaction with the Earth in this area does not need to be taken into account.}{1.20}

Note. There is no need to select a root sign at this point. Do not forget that within the framework of the model used, the interaction with the Earth in this area does not need to be taken into account.

Let us proceed to the analysis of the obtained expression for $v_1$. Next, the following relationships between quantities are used: $b/r_c \ll 1$ $a/r_m \ll 1$ $r_c/r_m \ll 1$

Let us proceed to the analysis of the obtained expression for $v_1$. The following relationships between quantities are used below:

\problembox{B4}{Show that the expression for $v_1$ obtained in paragraph B4, when expanded to first order in powers of $b/r_c$, is reduced to the form: $$v_1 \textbackslash\{\}approx \textbackslash\{\}dfrac\{v_m c\}\{s\} + B_0\textbackslash\{\}cdot \textbackslash\{\}dfrac\{b\}\{r_c\}$$ Выразите $B_0$ через $v_m$, $v_\text{II}^m$, $s$, $c$ и $s'$. Выбирая знак корня, покажите, что $B_0 > 0$.}{1.00}

$$v_1 \textbackslash\{\}approx \textbackslash\{\}dfrac\{v_m c\}\{s\} + B_0\textbackslash\{\}cdot \textbackslash\{\}dfrac\{b\}\{r_c\}$$

Выразите $B_0$ через $v_m$, $v_\text{II}^m$, $s$, $c$ и $s'$. Выбирая знак корня, покажите, что $B_0 > 0$.

Далее во всех пунктах задачи считайте, что дополнительно выполнено $\alpha \ll 1$.

\problembox{B5}{Show that the expression for $v_1$, up to first order, is: $$v_1 \textbackslash\{\}approx \textbackslash\{\}dfrac\{v_m\}\{\textbackslash\{\}text\{tg\}\textasciitilde\{\}\textbackslash\{\}theta\} + A \textbackslash\{\}alpha + B \textbackslash\{\}frac\{b\}\{r_c\}$$ Выразите $A$ и $B$ через $v_m$, $v_\text{II}^m$, $\theta$.}{0.60}

$$v_1 \textbackslash\{\}approx \textbackslash\{\}dfrac\{v_m\}\{\textbackslash\{\}text\{tg\}\textasciitilde\{\}\textbackslash\{\}theta\} + A \textbackslash\{\}alpha + B \textbackslash\{\}frac\{b\}\{r_c\}$$

Выразите $A$ и $B$ через $v_m$, $v_\text{II}^m$, $\theta$.

\problembox{B6}{Let's return to the Earth's reference system to find an expression for the angle $\alpha$ between the speed $v_1$ and the straight line Earth - Moon. Find an expression for $\alpha$ between the velocity $v_1$ and the Earth-Moon straight line up to first order in powers $r_c/r_m$ and $a/r_m$. Express your answer in terms of $a$, $r_c$, $r_m$, $v_1$, $v_0$ and $\theta$. Note. Do not forget that within the framework of the model used, interaction with the Moon in this area does not need to be taken into account.}{0.60}

Find an expression for $\alpha$ between the velocity $v_1$ and the Earth-Moon straight line up to first order in powers $r_c/r_m$ and $a/r_m$. Express your answer in terms of $a$, $r_c$, $r_m$, $v_1$, $v_0$ and $\theta$.

Note. Do not forget that within the framework of the model used, interaction with the Moon in this area does not need to be taken into account.

\problembox{B7}{Writing the ZSE in the FR of the Earth and neglecting the terms of order $a/r_m$ and $r_c/r_m$, express the speed $v_0$ in terms of $v_1$ and $v_\text{II}^e$. Note. Do not forget that within the framework of the model used, interaction with the Moon in this area does not need to be taken into account.}{0.40}

Note. Do not forget that within the framework of the model used, interaction with the Moon in this area does not need to be taken into account.

Let's return to the CO of the Moon again. Until now, we have never used the orbital symmetry condition anywhere. Let's use this condition and construct a hyperbola that is symmetrical with respect to the Earth-Moon straight line and passes through given points: the entry and exit points to and from the Moon's gravitational zone and the pericenter of the orbit.

\problembox{B8}{By writing the equation of a conic section, you get an equation from which you can find $\theta$, through $v_1$, $\alpha$, $v_m$, $G$, $M_m$, $b$ and $r_c$.}{1.00}

The resulting system of equations for $v_1$, $\alpha$ and $\theta$ can be approximately solved numerically, but we will not deal with this in this problem. The resulting speed $v_1$ will be an order of magnitude less than $v_\text{II}^e$. Then $v_0 = v_\text{II}^e$.

\problembox{B9}{Find $\beta_2$ - the ratio of the initial mass of the rocket to the final mass when accelerating from the speed $v_a$ - the orbital speed when moving in a circular orbit of radius $a$ to the speed $v_0$. Express your answer in terms of $u$, $v_\text{II}^e$ and $v_a$. Calculate the numerical value of $\beta_2$, take the quantities necessary for calculations from part A.}{0.40}

Calculate the numerical value of $\beta_2$, take the quantities necessary for calculations from part A.

\section*{Part C. Lunar landing (1.0 points)}

In this part we will look at the lunar landing. We will consider the case of a close approach to the Moon; for numerical calculations, consider $b = 1800$ km. Considering that $h = b - R_m \ll R_m$, we can assume with good accuracy that fuel is spent only on braking to zero speed at the point closest to the Moon.

\problembox{C1}{Considering $v_2 = 1.5~км/с$ (take the remaining necessary quantities from the previous parts of the problem), calculate numerically how many times the mass of the rocket will decrease during the landing process $\beta_3$.}{0.50}

\problembox{C2}{Calculate numerically what initial mass $m$ the rocket must have to deliver cargo weighing $40~т$ to the Moon.}{0.50}

\vspace{1em}
\noindent\textit{Source:} \url{https://pho.rs/p/5029}

\end{document}